% !TeX program = lualatex
\documentclass[12pt,aspectratio=169]{beamer}
\usepackage[utf8]{inputenc}
\usepackage[T1]{fontenc}
\usepackage{amsmath,amsfonts,amssymb}
\usepackage{graphicx}
\usepackage{xcolor}
\usepackage{hyperref}
\usepackage{anyfontsize}
\usepackage{lmodern}
\usepackage{longtable}
\usepackage{array}
\usepackage{booktabs}

% ====== EXTREME FONT REDUCTION ======
\renewcommand{\tiny}{\fontsize{3.5}{4.5}\selectfont}
\renewcommand{\scriptsize}{\fontsize{4.5}{5.5}\selectfont}
\renewcommand{\footnotesize}{\fontsize{5.5}{6.5}\selectfont}
\renewcommand{\small}{\fontsize{6.5}{7.5}\selectfont}
\renewcommand{\normalsize}{\fontsize{7.5}{8.5}\selectfont}
\renewcommand{\large}{\fontsize{8.5}{9.5}\selectfont}
\renewcommand{\Large}{\fontsize{9.5}{10.5}\selectfont}
\renewcommand{\LARGE}{\fontsize{10.5}{12}\selectfont}
\renewcommand{\huge}{\fontsize{11.5}{13.5}\selectfont}
\renewcommand{\Huge}{\fontsize{13}{16}\selectfont}

% ====== COLORS ======
\definecolor{redcorrect}{RGB}{200,30,30}
\definecolor{bluecorrect}{RGB}{30,80,200}
\definecolor{greencheck}{RGB}{0,150,50}
\definecolor{lightgray}{RGB}{245,245,245}
\definecolor{darkgray}{RGB}{40,40,40}
\definecolor{softyellow}{RGB}{255,248,200}
\definecolor{deepblue}{RGB}{20,60,150}
\definecolor{darkred}{RGB}{150,20,20}
\definecolor{orangewarning}{RGB}{255,140,0}
\definecolor{correctgreen}{RGB}{0,120,40}
\definecolor{trapcolor}{RGB}{180,80,0}
\definecolor{purpleconcept}{RGB}{120,50,180}
\definecolor{tealhard}{RGB}{0,120,120}

\setbeamercolor{title}{fg=white,bg=redcorrect}
\setbeamercolor{frametitle}{fg=white,bg=redcorrect}
\setbeamercolor{block title}{fg=white,bg=bluecorrect}
\setbeamercolor{block body}{fg=darkgray,bg=lightgray}
\setbeamercolor{normal text}{fg=darkgray}
\usecolortheme{default}

% ====== CUSTOM COMMANDS ======
\newcommand{\correct}[1]{\textcolor{greencheck}{\textbf{\checkmark}} #1}
\newcommand{\keypoint}[1]{\textcolor{deepblue}{\textbf{#1}}}
\newcommand{\hardconcept}[1]{\textcolor{darkred}{\textbf{#1}}}
\newcommand{\examtraps}[1]{\textcolor{trapcolor}{\textbf{⚠ TRAP:}} #1}
\newcommand{\recallbox}[1]{\fcolorbox{deepblue}{blue!5}{\parbox{0.88\textwidth}{\centering \textbf{REMEMBER:} #1}}}
\newcommand{\stepbox}[1]{%
	\begin{center}
		\fcolorbox{darkgray}{white}{%
			\parbox{0.92\textwidth}{\vspace{0.03cm} #1 \vspace{0.03cm}}%
		}
	\end{center}
}
\newcommand{\answerbox}[1]{\fcolorbox{correctgreen}{green!8}{\parbox{0.88\textwidth}{\centering \textcolor{correctgreen}{\textbf{\checkmark\ CORRECT:}} #1}}}
\newcommand{\wronganswer}[1]{\textcolor{redcorrect}{\textbf{✘}} #1}
\newcommand{\trapbox}[1]{\fcolorbox{trapcolor}{orange!8}{\parbox{0.88\textwidth}{\centering \textcolor{trapcolor}{\textbf{⚠ EXAM TRAP:}} #1}}}
\newcommand{\keyrule}[1]{\textcolor{tealhard}{\textbf{🔑}} #1}

\begin{document}
	
	% ================================================================
	% SLIDE 1: TITLE
	% ================================================================
	\begin{frame}
		\centering
		\vspace{0.3cm}
		{\fontsize{18}{22}\selectfont \textbf{CAMS DOMAIN C}}\\[0.1cm]
		{\fontsize{14}{17}\selectfont \textbf{BUILDING AN AML/AFC COMPLIANCE PROGRAM}}\\[0.15cm]
		{\fontsize{9}{11}\selectfont \textit{ALL 50+ — COMPLETE EXPANDED TUTORIAL}}\\[0.08cm]
		\rule{4cm}{0.5pt}\\[0.08cm]
		{\fontsize{6}{7}\selectfont MiCA · OFAC · Biometrics · MLAT · EWRA · DORA · TCSP · PEP EDD}\\[0.04cm]
		{\fontsize{6}{7}\selectfont KYC/CDD/EDD · Transaction Monitoring · Data Lineage · Clustering}\\[0.04cm]
		{\fontsize{6}{7}\selectfont Digital Onboarding · Deepfake Detection · AML Communications · PSP}\\[0.04cm]
		{\fontsize{6}{7}\selectfont Risk Assessments · Three Lines · Board Oversight · SAR/STR Reporting}
		\vspace{0.2cm}
	\end{frame}
	
	% ================================================================
	% SLIDE 2: AML POLICY MANDATORY ELEMENTS — THE SELECTION TRAP
	% ================================================================
	\begin{frame}
		\frametitle{\fontsize{11}{13}\selectfont \textbf{AML Policy — Mandatory Elements Trap}}
		\begin{block}{\fontsize{7}{9}\selectfont \textbf{The Hard Question — CAMS Exam Favorite}}
			\scriptsize
			Which of the following are key elements that AML/CFT policies must include? (Select Two.)
		\end{block}
		\vspace{0.02cm}
		\answerbox{\large \textbf{1. Assignment of compliance responsibilities. 2. Program review triggers based on regulatory or business events.}}
		\vspace{0.02cm}
		\begin{block}{\fontsize{7}{9}\selectfont \textbf{Natural Explanation — Why You Got This Wrong}}
			\scriptsize
			\begin{itemize}
				\item \wronganswer{Cash flow performance figures segmented by region} — \textbf{NOT} AML policy content. This is financial performance.
				\item \wronganswer{Marketing strategy} — \textbf{NOT} an AML policy element.
				\item \textbf{Assignment of compliance responsibilities:} Who does what? Must designate a Compliance Officer (MLRO in the UK) with clear authority and resources.
				\item \textbf{Program review triggers:} When to update the program? New regulations, significant business changes, new products/services, or changes in risk appetite.
				\item \textbf{BSA Requirement:} Written AML program must include: Internal controls, independent testing, designated compliance officer, training.
				\item \textbf{Key principle:} AML policies are \textbf{risk-based} — tailored to the institution's specific risks.
			\end{itemize}
		\end{block}
		\vspace{0.02cm}
		\recallbox{\scriptsize \textbf{Key Rule:} AML policy = \textbf{compliance responsibilities} + \textbf{review triggers} + internal controls + training + testing.}
	\end{frame}
	
	% ================================================================
	% SLIDE 3: TCSP — THE CONFIDENTIALITY AND OUTSOURCING TRAP
	% ================================================================
	\begin{frame}
		\frametitle{\fontsize{11}{13}\selectfont \textbf{TCSP — The Confidentiality and Outsourcing Trap}}
		\begin{block}{\fontsize{7}{9}\selectfont \textbf{The Scenario — CAMS Exam Favorite}}
			\scriptsize
			A trust and company service provider (TCSP) firm is reviewing whether to accept a client who wants to incorporate in a low-tax Caribbean jurisdiction. The potential client requests confidentiality protections for all shareholders. A connection is identified to industries marked high-risk in the NRA.
		\end{block}
		\vspace{0.02cm}
		\begin{block}{\fontsize{7}{9}\selectfont \textbf{The Hard Question}}
			\scriptsize
			What should the TCSP response be?
		\end{block}
		\vspace{0.02cm}
		\answerbox{\large \textbf{Decline immediately unless full beneficial owner information and risk justifications are disclosed and verified.}}
		\vspace{0.02cm}
		\begin{block}{\fontsize{7}{9}\selectfont \textbf{Natural Explanation — Why You Got This Wrong}}
			\scriptsize
			\begin{itemize}
				\item \wronganswer{Refer to offshore local counsel and detach from compliance responsibility} — \textbf{NO.} TCSP retains compliance responsibility. Cannot outsource AML obligations.
				\item \wronganswer{Accept with standard CDD} — \textbf{NO.} High-risk indicators require EDD or decline.
				\item \textbf{Confidentiality requests + high-risk industries + low-tax jurisdiction = significant red flags.}
				\item \textbf{FATF Rec 24/25:} BO identification is mandatory for legal persons and arrangements.
				\item \textbf{Offloading responsibility} to local counsel = \textbf{not} a valid risk mitigation strategy.
				\item \textbf{Key principle:} You can use local counsel, but you \textbf{cannot transfer liability}. You remain responsible.
			\end{itemize}
		\end{block}
		\vspace{0.02cm}
		\recallbox{\scriptsize \textbf{Key Rule:} TCSP cannot outsource AML responsibility. Decline unless BO info is disclosed and verified.}
	\end{frame}
	
	% ================================================================
	% SLIDE 4: AFC GOVERNANCE REPORTS — THE CONTENT TRAP
	% ================================================================
	\begin{frame}
		\frametitle{\fontsize{11}{13}\selectfont \textbf{AFC Governance Reports — The Content Trap}}
		\begin{block}{\fontsize{7}{9}\selectfont \textbf{The Hard Question — CAMS Exam Favorite}}
			\scriptsize
			Which of the following are appropriate contents of internal AFC reporting submitted to governance committees? (Select Three.)
		\end{block}
		\vspace{0.02cm}
		\answerbox{\large \textbf{1. Updates on regulatory interactions and supervisory reviews. 2. Assessment outcomes of control effectiveness on monitoring tools. 3. Investigation statuses of escalated SAR reviews.}}
		\vspace{0.02cm}
		\begin{block}{\fontsize{7}{9}\selectfont \textbf{Natural Explanation — Why You Got This Wrong}}
			\scriptsize
			\begin{itemize}
				\item \wronganswer{Cash flow performance figures segmented by region} — \textbf{NOT} a core AFC governance report component.
				\item \wronganswer{Sales targets by business unit} — \textbf{NOT} AFC governance content.
				\item \textbf{Regulatory interactions:} What did regulators say? Any findings? Any enforcement actions?
				\item \textbf{Control effectiveness:} Are our monitoring tools working? KPIs/KRIs on alerts, false positives, SARs.
				\item \textbf{SAR investigation statuses:} How many SARs filed? What's the disposition? Any law enforcement feedback?
				\item \textbf{Key principle:} Governance reports = \textbf{risk and compliance}, NOT financial performance.
			\end{itemize}
		\end{block}
		\vspace{0.02cm}
		\recallbox{\scriptsize \textbf{Key Rule:} AFC governance reports = \textbf{regulatory interactions, control effectiveness, SAR statuses}. NOT financial metrics.}
	\end{frame}
	
	% ================================================================
	% SLIDE 5: BOARD OVERSIGHT — THE OPERATIONAL INVOLVEMENT TRAP
	% ================================================================
	\begin{frame}
		\frametitle{\fontsize{11}{13}\selectfont \textbf{Board Oversight — The Operational Trap}}
		\begin{block}{\fontsize{7}{9}\selectfont \textbf{The Hard Question — CAMS Exam Favorite}}
			\scriptsize
			Which actions demonstrate effective board oversight of an AFC program? (Select Three.)
		\end{block}
		\vspace{0.02cm}
		\answerbox{\large \textbf{1. Receiving regular reports on financial crime risks and control effectiveness. 2. Reviewing whether senior management is addressing audit findings. 3. Challenging management when AFC performance metrics show signs of weakness.}}
		\vspace{0.02cm}
		\begin{block}{\fontsize{7}{9}\selectfont \textbf{Natural Explanation — Why You Got This Wrong}}
			\scriptsize
			\begin{itemize}
				\item \wronganswer{Participating in SAR investigations} — \textbf{Operational}, NOT governance.
				\item \wronganswer{Conducting transaction reviews} — \textbf{Operational}, NOT governance.
				\item \textbf{Board duties:}
				\begin{itemize}
					\item Receive reports (stay informed)
					\item Review audit responses (hold management accountable)
					\item Challenge management (ask tough questions)
				\end{itemize}
				\item \textbf{Key principle:} Board = \textbf{oversight and challenge}. NOT \textbf{doing the work}.
			\end{itemize}
		\end{block}
		\vspace{0.02cm}
		\recallbox{\scriptsize \textbf{Key Rule:} Board = \textbf{oversight}. NOT operational involvement.}
	\end{frame}
	
	% ================================================================
	% SLIDE 6: PASSIVE LIVENESS DETECTION — THE DISTINCTION TRAP
	% ================================================================
	\begin{frame}
		\frametitle{\fontsize{11}{13}\selectfont \textbf{Passive Liveness — The Distinction Trap}}
		\begin{block}{\fontsize{7}{9}\selectfont \textbf{The Hard Question — CAMS Exam Favorite}}
			\scriptsize
			What distinguishes passive liveness detection from active methods?
		\end{block}
		\vspace{0.02cm}
		\answerbox{\large \textbf{It uses minimal user interaction and analyzes subtle facial movement or texture patterns.}}
		\vspace{0.02cm}
		\begin{block}{\fontsize{7}{9}\selectfont \textbf{Natural Explanation — Why You Got This Wrong}}
			\scriptsize
			\begin{itemize}
				\item \wronganswer{It asks users to answer KYC security questions during video evaluation} — \textbf{NO.} That's KYC verification, NOT liveness detection.
				\item \textbf{Passive Liveness:}
				\begin{itemize}
					\item Uses \textbf{minimal user interaction}
					\item Analyzes subtle facial movement or texture patterns
					\item Micro-expressions, skin texture, depth perception
					\item Detects if using a photo or video (spoof detection)
				\end{itemize}
				\item \textbf{Active Liveness:}
				\begin{itemize}
					\item Requires explicit user actions: blinking, smiling, turning head
					\item More secure but more intrusive
					\item May have lower completion rates
				\end{itemize}
				\item \textbf{Key principle:} Passive is less intrusive; active is more secure.
			\end{itemize}
		\end{block}
		\vspace{0.02cm}
		\recallbox{\scriptsize \textbf{Key Rule:} Passive = \textbf{minimal interaction} + \textbf{subtle facial analysis}. Active = \textbf{user actions}.}
	\end{frame}
	
	% ================================================================
	% SLIDE 7: EWRA — THE DIGITAL CHANNEL TRAP
	% ================================================================
	\begin{frame}
		\frametitle{\fontsize{11}{13}\selectfont \textbf{EWRA — The Digital Channel Trap}}
		\begin{block}{\fontsize{7}{9}\selectfont \textbf{The Scenario — CAMS Exam Favorite}}
			\scriptsize
			A financial institution identifies that several of its high-risk jurisdiction customers conduct transactions through its mobile platform, bypassing in-branch monitoring measures.
		\end{block}
		\vspace{0.02cm}
		\begin{block}{\fontsize{7}{9}\selectfont \textbf{The Hard Question}}
			\scriptsize
			How should this be addressed as part of the enterprise-wide risk assessment (EWRA)?
		\end{block}
		\vspace{0.02cm}
		\answerbox{\large \textbf{Incorporate digital delivery channels into the EWRA to reassess risk exposure.}}
		\vspace{0.02cm}
		\begin{block}{\fontsize{7}{9}\selectfont \textbf{Natural Explanation — Why You Got This Wrong}}
			\scriptsize
			\begin{itemize}
				\item \wronganswer{Close all mobile accounts} — \textbf{De-risking}, NOT risk management.
				\item \wronganswer{Apply same controls as branches} — Digital channels have \textbf{different} risk profiles.
				\item \textbf{Mobile platforms often bypass branch monitoring} = risk blind spot.
				\item EWRA must be \textbf{updated} when new channels emerge.
				\item Digital-specific risk indicators must be added.
				\item \textbf{Key principle:} EWRA must be \textbf{dynamic} — updated when new channels, products, or customers emerge.
			\end{itemize}
		\end{block}
		\vspace{0.02cm}
		\recallbox{\scriptsize \textbf{Key Rule:} EWRA must \textbf{include digital channels}. Mobile = different risk profile.}
	\end{frame}
	
	% ================================================================
	% SLIDE 8: ENTITY RESOLUTION — THE BLOCKCHAIN TRAP
	% ================================================================
	\begin{frame}
		\frametitle{\fontsize{11}{13}\selectfont \textbf{Entity Resolution — The Blockchain Trap}}
		\begin{block}{\fontsize{7}{9}\selectfont \textbf{The Hard Question — CAMS Exam Favorite}}
			\scriptsize
			How does entity resolution improve blockchain analytics for on-chain behavior monitoring?
		\end{block}
		\vspace{0.02cm}
		\answerbox{\large \textbf{It links data from multiple sources to improve match accuracy.}}
		\vspace{0.02cm}
		\begin{block}{\fontsize{7}{9}\selectfont \textbf{Natural Explanation — Why You Got This Wrong}}
			\scriptsize
			\begin{itemize}
				\item \wronganswer{It replaces all manual reviews} — \textbf{NO.} It's a tool, not a replacement.
				\item \wronganswer{It eliminates false positives entirely} — \textbf{NO.} It reduces them but doesn't eliminate them.
				\item \textbf{Entity resolution:} Links on-chain wallet addresses to off-chain risk data.
				\item \textbf{Multiple sources:} Wallet addresses, KYC data, sanctions lists, PEP lists, adverse media.
				\item \textbf{Example:} A wallet address associated with ransomware is linked to an exchange account.
				\item \textbf{Key principle:} Entity resolution = \textbf{links data from multiple sources} to improve match accuracy.
			\end{itemize}
		\end{block}
		\vspace{0.02cm}
		\recallbox{\scriptsize \textbf{Key Rule:} Entity resolution = \textbf{links data from multiple sources} to improve accuracy.}
	\end{frame}
	
	% ================================================================
	% SLIDE 9: DATA COVERAGE AND GAP ASSESSMENT — THE CHALLENGES TRAP
	% ================================================================
	\begin{frame}
		\frametitle{\fontsize{11}{13}\selectfont \textbf{Data Coverage — The Challenges Trap}}
		\begin{block}{\fontsize{7}{9}\selectfont \textbf{The Hard Question — CAMS Exam Favorite}}
			\scriptsize
			Which of the following challenges can be identified through a data coverage and gap assessment? (Select Two.)
		\end{block}
		\vspace{0.02cm}
		\answerbox{\large \textbf{1. Missing identifiers in customer profiles. 2. Inconsistent use of country codes across systems.}}
		\vspace{0.02cm}
		\begin{block}{\fontsize{7}{9}\selectfont \textbf{Natural Explanation — Why You Got This Wrong}}
			\scriptsize
			\begin{itemize}
				\item \wronganswer{Misalignment between transaction monitoring scenarios and typologies} — This is a \textbf{scenario design issue}, not a data coverage gap.
				\item \textbf{Missing identifiers:} No DOB, no tax ID, no passport number — impacts screening against sanctions and PEP lists.
				\item \textbf{Inconsistent country codes:} "US" vs "USA" vs "840" vs "United States" — leads to false negatives and false positives in monitoring and reporting.
				\item \textbf{Key principle:} Data gaps and inconsistencies undermine the effectiveness of transaction monitoring, screening, and regulatory reporting.
			\end{itemize}
		\end{block}
		\vspace{0.02cm}
		\recallbox{\scriptsize \textbf{Key Rule:} Data gaps = \textbf{missing identifiers + inconsistent codes}.}
	\end{frame}
	
	% ================================================================
	% SLIDE 10: ROUND-TRIP TRANSACTIONS — THE ESCALATION TRAP
	% ================================================================
	\begin{frame}
		\frametitle{\fontsize{11}{13}\selectfont \textbf{Round-Trip Transactions — The Escalation Trap}}
		\begin{block}{\fontsize{7}{9}\selectfont \textbf{The Hard Question — CAMS Exam Favorite}}
			\scriptsize
			A Level 1 analyst reviews an alert triggered by a rules-based system identifying a "round trip" transaction. After reviewing recent account activity and confirming a nearly identical send and receive remittance within minutes, what should the analyst do next?
		\end{block}
		\vspace{0.02cm}
		\answerbox{\large \textbf{Escalate to Level 2 for deeper investigation.}}
		\vspace{0.02cm}
		\begin{block}{\fontsize{7}{9}\selectfont \textbf{Natural Explanation — Why You Got This Wrong}}
			\scriptsize
			\begin{itemize}
				\item \wronganswer{Close the alert as false positive} — \textbf{NO.} Round-trip transactions are suspicious.
				\item \wronganswer{Conduct full investigation at Level 1} — \textbf{NO.} Level 1 is for triage, not full investigation.
				\item \textbf{Level 1 analysts:} Initial triage and data gathering.
				\item \textbf{Level 2 analysts:} Deeper investigation, customer interview, SAR filing.
				\item \textbf{Round-trip:} Funds transferred out and returned to the same account — classic layering technique.
				\item \textbf{Key principle:} Escalate when the pattern is confirmed.
			\end{itemize}
		\end{block}
		\vspace{0.02cm}
		\recallbox{\scriptsize \textbf{Key Rule:} Round-trip = \textbf{escalate to Level 2}. Level 1 = triage.}
	\end{frame}
	
	% ================================================================
	% SLIDE 11: AI TRANSITION — THE DATA/PEOPLE/EXPLAINABILITY TRAP
	% ================================================================
	\begin{frame}
		\frametitle{\fontsize{11}{13}\selectfont \textbf{AI Transition — The Data/People/Explainability Trap}}
		\begin{block}{\fontsize{7}{9}\selectfont \textbf{The Hard Question — CAMS Exam Favorite}}
			\scriptsize
			Which actions are most important when transitioning to AI-based compliance tools? (Select Three.)
		\end{block}
		\vspace{0.02cm}
		\answerbox{\large \textbf{1. Ensuring data quality for training algorithms. 2. Managing change in employee workflows. 3. Verifying transparency and explainability in AI models.}}
		\vspace{0.02cm}
		\begin{block}{\fontsize{7}{9}\selectfont \textbf{Natural Explanation — Why You Got This Wrong}}
			\scriptsize
			\begin{itemize}
				\item \wronganswer{Selecting third-party global audit providers} — This is \textbf{less critical} than data, people, and explainability.
				\item \textbf{Data quality:} "Garbage in, garbage out" — AI is only as good as its training data.
				\item \textbf{Employee workflows:} AI changes how analysts work — must provide adequate training.
				\item \textbf{Explainability:} Regulators require \textbf{explainable AI} — must justify why an AI flagged something.
				\item \textbf{Key principle:} AI transition = \textbf{data + people + explainability}.
			\end{itemize}
		\end{block}
		\vspace{0.02cm}
		\recallbox{\scriptsize \textbf{Key Rule:} AI transition = \textbf{data + people + explainability}.}
	\end{frame}
	
	% ================================================================
	% SLIDE 12: CARD-BASED PRODUCT RISK FACTORS — THE SELECTION TRAP
	% ================================================================
	\begin{frame}
		\frametitle{\fontsize{11}{13}\selectfont \textbf{Card-Based Products — The Selection Trap}}
		\begin{block}{\fontsize{7}{9}\selectfont \textbf{The Hard Question — CAMS Exam Favorite}}
			\scriptsize
			Which of the following represent common risk factors associated with card-based products? (Select Three.)
		\end{block}
		\vspace{0.02cm}
		\answerbox{\large \textbf{1. Anonymity through prepaid cards. 2. High volumes of rapid transactions. 3. Use by customers in high-risk sectors.}}
		\vspace{0.02cm}
		\begin{block}{\fontsize{7}{9}\selectfont \textbf{Natural Explanation — Why You Got This Wrong}}
			\scriptsize
			\begin{itemize}
				\item \wronganswer{Low transaction limits} — Low limits \textbf{reduce} risk, not increase it.
				\item \textbf{Anonymity through prepaid cards:} Prepaid cards can be purchased with cash and used without identification.
				\item \textbf{High volumes of rapid transactions:} Card transactions are processed instantly in high volumes — structuring (smurfing) can occur.
				\item \textbf{Use by customers in high-risk sectors:} Merchants in crypto, gambling, or cash-intensive businesses generate high-risk card transactions.
				\item \textbf{Key principle:} Card risk = \textbf{anonymity + speed + sector}.
			\end{itemize}
		\end{block}
		\vspace{0.02cm}
		\recallbox{\scriptsize \textbf{Key Rule:} Card risk = \textbf{anonymity + high volume + high-risk sectors}.}
	\end{frame}
	
	% ================================================================
	% SLIDE 13: PREPAID CARDS — THE VULNERABILITY TRAP
	% ================================================================
	\begin{frame}
		\frametitle{\fontsize{11}{13}\selectfont \textbf{Prepaid Cards — The Vulnerability Trap}}
		\begin{block}{\fontsize{7}{9}\selectfont \textbf{The Hard Question — CAMS Exam Favorite}}
			\scriptsize
			Why are prepaid cards particularly vulnerable to financial crime risks?
		\end{block}
		\vspace{0.02cm}
		\answerbox{\large \textbf{They can be used anonymously, especially when reloadable or unregistered.}}
		\vspace{0.02cm}
		\begin{block}{\fontsize{7}{9}\selectfont \textbf{Natural Explanation — Why You Got This Wrong}}
			\scriptsize
			\begin{itemize}
				\item \wronganswer{They always require full KYC} — \textbf{NO.} Many prepaid cards do NOT require KYC.
				\item \wronganswer{They have high transaction limits} — \textbf{NO.} The risk is anonymity, not limits.
				\item \textbf{Unregistered cards:} Can be purchased with cash at retail outlets; no customer identification required.
				\item \textbf{Reloadable cards:} Can be topped up multiple times, potentially with illicit funds.
				\item \textbf{Cross-border use:} Can be used internationally, bypassing local AML controls.
				\item \textbf{Key principle:} Prepaid cards = \textbf{anonymity + reloadability + cross-border}.
			\end{itemize}
		\end{block}
		\vspace{0.02cm}
		\recallbox{\scriptsize \textbf{Key Rule:} Prepaid cards = \textbf{anonymous + reloadable + cross-border}.}
	\end{frame}
	
	% ================================================================
	% SLIDE 14: CASINO AML CONTROLS — THE SELECTION TRAP
	% ================================================================
	\begin{frame}
		\frametitle{\fontsize{11}{13}\selectfont \textbf{Casino Controls — The Selection Trap}}
		\begin{block}{\fontsize{7}{9}\selectfont \textbf{The Hard Question — CAMS Exam Favorite}}
			\scriptsize
			Which control measures help mitigate AML risks in casinos? (Select Two.)
		\end{block}
		\vspace{0.02cm}
		\answerbox{\large \textbf{1. Enforcing chip redemption policies tied to minimum gaming thresholds. 2. Tracking player activity through registered loyalty programs.}}
		\vspace{0.02cm}
		\begin{block}{\fontsize{7}{9}\selectfont \textbf{Natural Explanation — Why You Got This Wrong}}
			\scriptsize
			\begin{itemize}
				\item \wronganswer{Monitoring unstructured table gameplay for win/loss patterns} — This is a \textbf{monitoring technique}, not a formal control measure.
				\item \textbf{Chip redemption policies:} Prevents "chip washing" — exchanging chips for cash without meaningful gaming.
				\item \textbf{Thresholds:} US \$10,000 triggers reporting.
				\item \textbf{Loyalty programs:} Provide rich data on player behavior (wins, losses, session durations).
				\item \textbf{Key principle:} Casino controls = \textbf{chip redemption + loyalty program tracking}.
			\end{itemize}
		\end{block}
		\vspace{0.02cm}
		\recallbox{\scriptsize \textbf{Key Rule:} Casino controls = \textbf{chip redemption + loyalty program tracking}.}
	\end{frame}
	
	% ================================================================
	% SLIDE 15: E-COMMERCE TRANSACTION MONITORING — THE VOLUME TRAP
	% ================================================================
	\begin{frame}
		\frametitle{\fontsize{11}{13}\selectfont \textbf{E-Commerce Monitoring — The Volume Trap}}
		\begin{block}{\fontsize{7}{9}\selectfont \textbf{The Hard Question — CAMS Exam Favorite}}
			\scriptsize
			Why is transaction monitoring particularly challenging in an e-commerce risk context?
		\end{block}
		\vspace{0.02cm}
		\answerbox{\large \textbf{High transaction volumes and decentralized payments make patterns harder to detect.}}
		\vspace{0.02cm}
		\begin{block}{\fontsize{7}{9}\selectfont \textbf{Natural Explanation — Why You Got This Wrong}}
			\scriptsize
			\begin{itemize}
				\item \wronganswer{Transactions are always small and therefore low-risk} — \textbf{NO.} Small transactions can still be suspicious.
				\item \textbf{High volume:} Millions of transactions daily — traditional rules-based monitoring generates too many false positives.
				\item \textbf{Decentralized payments:} Multiple payment gateways, cryptocurrencies, digital wallets.
				\item \textbf{Consequence:} Patterns are harder to detect.
				\item \textbf{Mitigation:} Use AI/ML models to detect anomalies; segment data by merchant, payment method, and product.
				\item \textbf{Key principle:} E-commerce challenge = \textbf{high volume + decentralized payments}.
			\end{itemize}
		\end{block}
		\vspace{0.02cm}
		\recallbox{\scriptsize \textbf{Key Rule:} E-commerce challenge = \textbf{high volume + decentralized payments}.}
	\end{frame}
	
	% ================================================================
	% SLIDE 16: PSP — THE CUSTOMIZATION TRAP
	% ================================================================
	\begin{frame}
		\frametitle{\fontsize{11}{13}\selectfont \textbf{PSP — The Customization Trap}}
		\begin{block}{\fontsize{7}{9}\selectfont \textbf{The Scenario — CAMS Exam Favorite}}
			\scriptsize
			A PSP is planning to offer new services targeting small merchants in multiple countries, some of which lack mature AML frameworks.
		\end{block}
		\vspace{0.02cm}
		\begin{block}{\fontsize{7}{9}\selectfont \textbf{The Hard Question}}
			\scriptsize
			Which strategy is most appropriate to maintain compliance and mitigate risk?
		\end{block}
		\vspace{0.02cm}
		\answerbox{\large \textbf{Customize AML policy and onboarding procedures per jurisdiction and product.}}
		\vspace{0.02cm}
		\begin{block}{\fontsize{7}{9}\selectfont \textbf{Natural Explanation — Why You Got This Wrong}}
			\scriptsize
			\begin{itemize}
				\item \wronganswer{Apply home-country AML standards across all jurisdictions} — \textbf{NO.} One-size-fits-all fails.
				\item \textbf{Jurisdiction-specific:} High-risk jurisdictions = EDD, transaction limits, or restricted services.
				\item \textbf{Product-specific:} Low-value vs high-value products have different risk profiles.
				\item \textbf{Risk-based approach:} Apply proportionate measures based on risk.
				\item \textbf{Key principle:} PSPs must \textbf{customize} AML per jurisdiction and product.
			\end{itemize}
		\end{block}
		\vspace{0.02cm}
		\recallbox{\scriptsize \textbf{Key Rule:} PSPs = \textbf{customize per jurisdiction and product}.}
	\end{frame}
	
	% ================================================================
	% SLIDE 17: HIGH-VALUE ASSETS — THE CRYPTO RED FLAG TRAP
	% ================================================================
	\begin{frame}
		\frametitle{\fontsize{11}{13}\selectfont \textbf{High-Value Assets — The Crypto Red Flag Trap}}
		\begin{block}{\fontsize{7}{9}\selectfont \textbf{The Hard Question — CAMS Exam Favorite}}
			\scriptsize
			Which red flag is most commonly associated with abuse of high-value assets for laundering purposes?
		\end{block}
		\vspace{0.02cm}
		\answerbox{\large \textbf{Making payments for high-value items in cryptocurrency.}}
		\vspace{0.02cm}
		\begin{block}{\fontsize{7}{9}\selectfont \textbf{Natural Explanation — Why You Got This Wrong}}
			\scriptsize
			\begin{itemize}
				\item \wronganswer{Purchasing assets with complex financing plans spanning several years} — This \textbf{can} be a red flag, but crypto is \textbf{more directly} associated.
				\item \textbf{Cryptocurrency offers pseudonymity}, making it attractive for illicit payments.
				\item \textbf{High-value items:} Real estate, art, luxury cars purchased with crypto = suspicion.
				\item \textbf{Challenge:} Volatility and price fluctuations complicate source-of-funds verification.
				\item \textbf{Key principle:} Crypto payments for high-value assets = \textbf{primary red flag}.
			\end{itemize}
		\end{block}
		\vspace{0.02cm}
		\recallbox{\scriptsize \textbf{Key Rule:} Crypto payments for high-value assets = \textbf{primary red flag}.}
	\end{frame}
	
	% ================================================================
	% SLIDE 18: AML COMMUNICATIONS — THE SCRIPT TRAP
	% ================================================================
	\begin{frame}
		\frametitle{\fontsize{11}{13}\selectfont \textbf{AML Communications — The Script Trap}}
		\begin{block}{\fontsize{7}{9}\selectfont \textbf{The Hard Question — CAMS Exam Favorite}}
			\scriptsize
			What should AML-sensitive communications materials emphasize to reduce risk? (Select Two.)
		\end{block}
		\vspace{0.02cm}
		\answerbox{\large \textbf{1. Avoid mentioning any internal reviews, even if prompted by the customer. 2. Use pre-approved scripts when addressing service disruptions caused by internal reviews.}}
		\vspace{0.02cm}
		\begin{block}{\fontsize{7}{9}\selectfont \textbf{Natural Explanation — Why You Got This Wrong}}
			\scriptsize
			\begin{itemize}
				\item \wronganswer{Refer customers to the AML department for financial advice} — \textbf{NO.} AML department doesn't give financial advice.
				\item \textbf{Avoid mentioning internal reviews:} This prevents tipping off.
				\item \textbf{Pre-approved scripts:} Ensures consistent, compliant messaging.
				\item \textbf{Tipping off risk:} Even hinting at an investigation can be tipping off.
				\item \textbf{Key principle:} Customer-facing staff must \textbf{not} reveal any AML-related investigations.
			\end{itemize}
		\end{block}
		\vspace{0.02cm}
		\recallbox{\scriptsize \textbf{Key Rule:} AML communications = \textbf{avoid reviews + use pre-approved scripts}.}
	\end{frame}
	
	% ================================================================
	% SLIDE 19: DEEPFAKE DETECTION — THE LIVENESS TRAP
	% ================================================================
	\begin{frame}
		\frametitle{\fontsize{11}{13}\selectfont \textbf{Deepfake Detection — The Liveness Trap}}
		\begin{block}{\fontsize{7}{9}\selectfont \textbf{The Scenario — CAMS Exam Favorite}}
			\scriptsize
			A fraudster used a high-resolution image and a deepfake video to bypass facial recognition verification.
		\end{block}
		\vspace{0.02cm}
		\begin{block}{\fontsize{7}{9}\selectfont \textbf{The Hard Question}}
			\scriptsize
			What tools could have prevented this fraud attempt? (Select Two.)
		\end{block}
		\vspace{0.02cm}
		\answerbox{\large \textbf{1. Use of software that can detect inconsistencies in light or skin texture. 2. Prompts for the customer to blink, smile, or repeat a phrase.}}
		\vspace{0.02cm}
		\begin{block}{\fontsize{7}{9}\selectfont \textbf{Natural Explanation — Why You Got This Wrong}}
			\scriptsize
			\begin{itemize}
				\item \wronganswer{Real-time fraud risk scoring algorithms using behavioral analytics} — This is \textbf{behavioral analytics}, not deepfake detection.
				\item \textbf{Texture analysis:} Deepfakes often have inconsistencies in light, skin texture, and shadows.
				\item \textbf{Active liveness prompts:} Blink, smile, repeat a phrase — requires real-time action.
				\item \textbf{Passive liveness:} Analyzes subtle facial movements without user action.
				\item \textbf{Key principle:} Deepfake detection = \textbf{texture analysis + active liveness prompts}.
			\end{itemize}
		\end{block}
		\vspace{0.02cm}
		\recallbox{\scriptsize \textbf{Key Rule:} Deepfake detection = \textbf{texture analysis + active liveness prompts}.}
	\end{frame}
	
	% ================================================================
	% SLIDE 20: CLUSTERING — THE BEHAVIORAL TRAP
	% ================================================================
	\begin{frame}
		\frametitle{\fontsize{11}{13}\selectfont \textbf{Clustering — The Behavioral Trap}}
		\begin{block}{\fontsize{7}{9}\selectfont \textbf{The Hard Question — CAMS Exam Favorite}}
			\scriptsize
			Which are characteristics of clustering in compliance analytics? (Select Two.)
		\end{block}
		\vspace{0.02cm}
		\answerbox{\large \textbf{1. It groups entities by shared behavioral characteristics. 2. It identifies connections across multiple transaction types.}}
		\vspace{0.02cm}
		\begin{block}{\fontsize{7}{9}\selectfont \textbf{Natural Explanation — Why You Got This Wrong}}
			\scriptsize
			\begin{itemize}
				\item \wronganswer{It excludes outliers to reduce financial crime risk} — \textbf{NO.} Outliers are often the \textbf{most suspicious}.
				\item \textbf{Clustering:} Groups entities that exhibit similar behaviors.
				\item \textbf{Purpose:} Detect networks of criminal activity (money mule networks, organized crime).
				\item \textbf{Connections across transaction types:} Links separate accounts that exhibit similar patterns.
				\item \textbf{Key principle:} Clustering = \textbf{shared behavior + connections}. NOT excluding outliers.
			\end{itemize}
		\end{block}
		\vspace{0.02cm}
		\recallbox{\scriptsize \textbf{Key Rule:} Clustering = \textbf{shared behaviors + connections}.}
	\end{frame}
	
	% ================================================================
	% SLIDE 21: DATA LINEAGE — THE TRANSPARENCY TRAP
	% ================================================================
	\begin{frame}
		\frametitle{\fontsize{11}{13}\selectfont \textbf{Data Lineage — The Transparency Trap}}
		\begin{block}{\fontsize{7}{9}\selectfont \textbf{The Hard Question — CAMS Exam Favorite}}
			\scriptsize
			Which support transparency and auditability through data lineage? (Select Two.)
		\end{block}
		\vspace{0.02cm}
		\answerbox{\large \textbf{1. Maintaining history of data transformations. 2. Capturing source of each data field.}}
		\vspace{0.02cm}
		\begin{block}{\fontsize{7}{9}\selectfont \textbf{Natural Explanation — Why You Got This Wrong}}
			\scriptsize
			\begin{itemize}
				\item \wronganswer{Archiving analyst notes from alert reviews} — This is \textbf{case documentation}, NOT data lineage.
				\item \wronganswer{Deleting outdated data annually} — \textbf{Undermines} auditability.
				\item \textbf{Data lineage:} Tracking the \textbf{flow of data} from source to destination.
				\item \textbf{History of transformations:} Shows how data was modified.
				\item \textbf{Source of each field:} Shows where data originated.
				\item \textbf{Key principle:} Data lineage = \textbf{transformations + source}. NOT case notes.
			\end{itemize}
		\end{block}
		\vspace{0.02cm}
		\recallbox{\scriptsize \textbf{Key Rule:} Data lineage = \textbf{transformations history + source of each field}.}
	\end{frame}
	
	% ================================================================
	% SLIDE 22: DATA VALIDATION — THE OBJECTIVES TRAP
	% ================================================================
	\begin{frame}
		\frametitle{\fontsize{11}{13}\selectfont \textbf{Data Validation — The Objectives Trap}}
		\begin{block}{\fontsize{7}{9}\selectfont \textbf{The Hard Question — CAMS Exam Favorite}}
			\scriptsize
			Which are typical objectives of data validation and testing in AFC systems? (Select Three.)
		\end{block}
		\vspace{0.02cm}
		\answerbox{\large \textbf{1. Ensuring data inputs reflect actual customer behavior. 2. Verifying automated controls function correctly. 3. Testing whether controls generate alerts under appropriate risk conditions.}}
		\vspace{0.02cm}
		\begin{block}{\fontsize{7}{9}\selectfont \textbf{Natural Explanation — Why You Got This Wrong}}
			\scriptsize
			\begin{itemize}
				\item \wronganswer{Confirming reports are generated in a format preferred by third-party vendors} — \textbf{NOT} a validation objective.
				\item \textbf{Data inputs reflect behavior:} Test that data from source systems accurately feeds into AML systems.
				\item \textbf{Controls functioning:} Validate that transaction monitoring rules fire alerts correctly.
				\item \textbf{Alerts under risk conditions:} Use known test cases to validate system logic.
				\item \textbf{Key principle:} Validation = \textbf{accuracy + functionality + appropriate alerts}.
			\end{itemize}
		\end{block}
		\vspace{0.02cm}
		\recallbox{\scriptsize \textbf{Key Rule:} Data validation = \textbf{accuracy + functionality + appropriate alerts}.}
	\end{frame}
	
	% ================================================================
	% SLIDE 23: DIGITAL ONBOARDING — THE BENEFITS TRAP
	% ================================================================
	\begin{frame}
		\frametitle{\fontsize{11}{13}\selectfont \textbf{Digital Onboarding — The Benefits Trap}}
		\begin{block}{\fontsize{7}{9}\selectfont \textbf{The Hard Question — CAMS Exam Favorite}}
			\scriptsize
			Which benefits are associated with digital onboarding technology according to FATF? (Select Two.)
		\end{block}
		\vspace{0.02cm}
		\answerbox{\large \textbf{1. Lower cost and increased auditability. 2. Greater accountability in compliance processes.}}
		\vspace{0.02cm}
		\begin{block}{\fontsize{7}{9}\selectfont \textbf{Natural Explanation — Why You Got This Wrong}}
			\scriptsize
			\begin{itemize}
				\item \wronganswer{Elimination of human review for high-risk customers} — \textbf{NO.} High-risk still require human review.
				\item \textbf{Lower cost:} Automation reduces manual effort.
				\item \textbf{Increased auditability:} Digital records are easier to track and audit.
				\item \textbf{Greater accountability:} Clear audit trails of who did what.
				\item \textbf{Key principle:} Digital onboarding = \textbf{lower cost + auditability + accountability}. NOT eliminating human review.
			\end{itemize}
		\end{block}
		\vspace{0.02cm}
		\recallbox{\scriptsize \textbf{Key Rule:} Digital onboarding = \textbf{lower cost + auditability + accountability}.}
	\end{frame}
	
	% ================================================================
	% SLIDE 24: PERPETUAL KYC — THE TRIGGER-BASED TRAP
	% ================================================================
	\begin{frame}
		\frametitle{\fontsize{11}{13}\selectfont \textbf{Perpetual KYC — The Trigger-Based Trap}}
		\begin{block}{\fontsize{7}{9}\selectfont \textbf{The Hard Question — CAMS Exam Favorite}}
			\scriptsize
			Which feature most effectively distinguishes perpetual KYC from traditional periodic reviews?
		\end{block}
		\vspace{0.02cm}
		\answerbox{\large \textbf{Trigger-based updates rather than fixed calendar reviews.}}
		\vspace{0.02cm}
		\begin{block}{\fontsize{7}{9}\selectfont \textbf{Natural Explanation — Why You Got This Wrong}}
			\scriptsize
			\begin{itemize}
				\item \wronganswer{It eliminates all manual review} — \textbf{NO.} Human review is still required.
				\item \textbf{Perpetual KYC (pKYC):} Continuous, event-driven KYC refresh.
				\item \textbf{Triggers:}
				\begin{itemize}
					\item Changes in PEP status
					\item Adverse media hits
					\item Unusual transaction patterns
					\item Jurisdictional risk changes
				\end{itemize}
				\item \textbf{Benefits:} Reduces customer friction, improves accuracy, ensures timely EDD.
				\item \textbf{Key principle:} pKYC = \textbf{trigger-based}, not calendar-based.
			\end{itemize}
		\end{block}
		\vspace{0.02cm}
		\recallbox{\scriptsize \textbf{Key Rule:} pKYC = \textbf{trigger-based updates}, not fixed calendar reviews.}
	\end{frame}
	
	% ================================================================
	% SLIDE 25: FINAL EXAM TIPS — DOMAIN C COMPLETE
	% ================================================================
	\begin{frame}
		\frametitle{\fontsize{11}{13}\selectfont \textbf{Domain C — Complete Exam Traps Summary}}
		\scriptsize
		\begin{block}{\fontsize{7}{9}\selectfont \textbf{All Traps — 50+ Images Covered}}
			\scriptsize
			\begin{enumerate}
				\item \textbf{AML Policy:} Assignment of compliance responsibilities + program review triggers.
				\item \textbf{TCSP:} Cannot outsource AML responsibility. Decline unless BO disclosed.
				\item \textbf{AFC Governance Reports:} Regulatory interactions + control effectiveness + SAR statuses.
				\item \textbf{Board Oversight:} Oversight and challenge. NOT operational involvement.
				\item \textbf{Passive Liveness:} Minimal interaction + subtle facial analysis.
				\item \textbf{EWRA:} Must include digital channels. Update when risk changes.
				\item \textbf{Entity Resolution:} Links data from multiple sources to improve accuracy.
				\item \textbf{Data Coverage:} Missing identifiers + inconsistent country codes.
				\item \textbf{Round-Trip:} Escalate to Level 2 for deeper investigation.
				\item \textbf{AI Transition:} Data + people + explainability.
				\item \textbf{Card Risk:} Anonymity + high volume + high-risk sectors.
				\item \textbf{Prepaid Cards:} Anonymous + reloadable + cross-border.
				\item \textbf{Casino Controls:} Chip redemption + loyalty program tracking.
				\item \textbf{E-Commerce:} High volume + decentralized payments.
				\item \textbf{PSP:} Customize per jurisdiction and product.
				\item \textbf{High-Value Assets:} Crypto payments = primary red flag.
				\item \textbf{AML Communications:} Avoid reviews + use pre-approved scripts.
				\item \textbf{Deepfake Detection:} Texture analysis + active liveness prompts.
				\item \textbf{Clustering:} Shared behaviors + connections across transaction types.
				\item \textbf{Data Lineage:} Transformations history + source of each field.
				\item \textbf{Data Validation:} Accuracy + functionality + appropriate alerts.
				\item \textbf{Digital Onboarding:} Lower cost + auditability + accountability.
				\item \textbf{Perpetual KYC:} Trigger-based updates, not fixed calendar reviews.
				\item \textbf{Risk-Based Authentication:} Prompt only when behavior deviates.
				\item \textbf{MLAT:} Challenge = admissibility, not just sharing.
			\end{enumerate}
		\end{block}
		\vspace{0.02cm}
		\recallbox{\scriptsize \textbf{Final Rule:} Understand the \textbf{why} behind each concept. The exam tests \textbf{application}, not memorization.}
	\end{frame}
	
	% ================================================================
	% END SLIDE
	% ================================================================
	\begin{frame}
		\centering
		\vspace{1.5cm}
		{\fontsize{18}{22}\selectfont \textbf{Domain C — Complete}}\\[0.2cm]
		{\fontsize{11}{13}\selectfont All 50+  — Full Expanded Tutorial}\\[0.1cm]
		\rule{4cm}{0.5pt}\\[0.1cm]
		{\fontsize{8}{10}\selectfont MiCA · OFAC · Biometrics · MLAT · EWRA · DORA · TCSP · PEP EDD}\\[0.05cm]
		{\fontsize{8}{10}\selectfont KYC/CDD/EDD · Transaction Monitoring · Data Lineage · Clustering}\\[0.05cm]
		{\fontsize{8}{10}\selectfont Digital Onboarding · Deepfake Detection · AML Communications · PSP}\\[0.05cm]
		{\fontsize{8}{10}\selectfont Risk Assessments · Three Lines · Board Oversight · SAR/STR Reporting}\\[0.1cm]
		\rule{4cm}{0.5pt}\\[0.1cm]
		{\fontsize{8}{10}\selectfont \textit{All traps explained. All concepts covered. Ready for the exam.}}
		\vspace{0.5cm}
	\end{frame}
	
\end{document}